\subsection{Rotación Mínima, Booth (huellas duplicadas por rotación)}

\graybold{Preguntas guía:}

\begin{itemize}
\item ¿Dos strings se consideran "iguales" si uno es una rotación cíclica del otro?
\item ¿Necesito una forma canónica única para comparar/deduplicar strings bajo rotación, en tiempo lineal (no \bigO{n^2} probando todas las rotaciones)?
\item ¿El string puede ser largo (miles de caracteres) y hay muchos para procesar?
\end{itemize}

\graybold{Utilidad:} Encontrar la rotación lexicográficamente menor de un string en \bigO{n}, como forma canónica para comparar o deduplicar strings equivalentes bajo rotación cíclica. Evita el enfoque ingenuo de generar y comparar las $n$ rotaciones (\bigO{n^2}).

\graybold{Complejidad:} \bigO{n}.

\graybold{Fundamento teórico:} El algoritmo de Booth adapta la función de fallo (como en KMP) sobre el string duplicado ($s+s$). Mantiene un índice candidato $k$ a la rotación mínima y, al procesar cada carácter, usa la función de fallo para saltar comparaciones ya conocidas, actualizando $k$ solo cuando encuentra evidencia de que otra posición produce una rotación menor. El resultado equivale a comparar todas las rotaciones en un solo recorrido lineal.

\graybold{Ejercicio aplicado}

Dado un conjunto de huellas dactilares codificadas como strings (equivalentes bajo rotación cíclica), eliminar duplicados representando cada grupo por su rotación canónica (la lexicográficamente menor).

\graybold{I/O:} Hasta 100 strings de hasta 5000 caracteres por caso, múltiples casos hasta el centinela "0" → lectura total con \texttt{sys.stdin.read().split()} (patrón 2). Verificado: \aprox{}1.4s para 5 casos × 100 strings × 5000 caracteres, y correcto contra fuerza bruta en 3000 strings aleatorios.

\graybold{Código solución}

\begin{lstlisting}[language=Python]
import sys

def least_rotation(s):
    # indice de inicio de la rotacion lexicograficamente menor (algoritmo de Booth)
    ss = s + s
    n = len(ss)
    fail = [-1] * n
    k = 0
    for j in range(1, n):
        sj = ss[j]
        i = fail[j - k - 1]
        while i != -1 and sj != ss[k + i + 1]:
            if sj < ss[k + i + 1]:
                k = j - i - 1
            i = fail[i]
        if sj != ss[k + i + 1]:
            if sj < ss[k + i + 1]:
                k = j
            fail[j - k] = -1
        else:
            fail[j - k] = i + 1
    return k

def canonical(s):
    k = least_rotation(s)
    return s[k:] + s[:k]

def solve():
    data = sys.stdin.read().split()
    idx = 0
    out = []
    while idx < len(data):
        n = int(data[idx]); idx += 1
        if n == 0:
            break
        seen = set()
        for _ in range(n):
            m = int(data[idx]); f = data[idx + 1]; idx += 2
            c = canonical(f)
            if c not in seen:            # dedup preservando orden de aparicion
                seen.add(c)
                out.append(c)
    sys.stdout.write("\n".join(out) + "\n")

solve()
\end{lstlisting}
